% ===================================================================
%  TAULA N.º 15 — El Mercat. — Els Aliments.
%  Traduction 2026 d'apres texto/io, controlee sur texto/fr et
%  texto/es. Consigne : voir texto/ca/10-taula-01.tex.
% ===================================================================

%%K t15-apar-1 apar
\VUsaut{4.0mm}
\VUtitre{15.0pt}{60}{TAULA N.º 15}
\VUsaut{3.0mm}

%%K t15-tit-1 sub
\VUpk{11.4pt}{40}{El Mercat. --- Els Aliments.}
\VUsaut{3.0mm}

%%K t15-01-1 p
1. --- La major part dels comerciants que ens proveeixen dels aliments necessaris són reunits
sota la \VUgras{nau} \textsuperscript{(1)} del \VUgras{mercat cobert}, o als carrers veïns.

%%K t15-02-1 p
2. --- El \VUgras{cansalader} \textsuperscript{(2)}, que ven \VUgras{pernil}
\textsuperscript{(3)}, \VUgras{botifarra negra} \textsuperscript{(4)},
\VUgras{llonganissa} \textsuperscript{(5)}, \VUgras{bull}
\textsuperscript{(6)} i \VUgras{cansalada} \textsuperscript{(7)}, és vora la
\VUgras{venedora de mantega i formatge} \textsuperscript{(8)}, la parada de la qual
és coberta de \VUgras{formatges de bola} \textsuperscript{(9)} de crosta vermella,
\VUgras{camemberts} \textsuperscript{(10)}, grans \VUgras{rodes de gruyère}
\textsuperscript{(11)} i \VUgras{pans de mantega} \textsuperscript{(12)} fresca.

%%K t15-03-1 p
3. --- Davant d'ella, una \VUgras{verdulaire} \textsuperscript{(13)} ofereix als seus clients
bells \VUgras{tomàquets} \textsuperscript{(14)}, \VUgras{colsiflors} \textsuperscript{(15)},
\VUgras{enciams} \textsuperscript{(16)}, \VUgras{cols} \textsuperscript{(17)},
\VUgras{carbassons} \textsuperscript{(19)},
\VUgras{melons} \textsuperscript{(18)} i fins i tot \VUgras{bolets}
\textsuperscript{(20)}. Però un \VUgras{cerveser}, que ha aturat just davant de la seva parada
el seu \VUgras{carro de repartiment} \textsuperscript{(21)} carregat de
\VUgras{barrils de cervesa} \textsuperscript{(22)}, impedeix que els seus clients
s'acostin. Un hortolà li porta un cistell de \VUgras{carxofes} \textsuperscript{(23)}.

%%K t15-tit-2 sub
\VUsaut{4.0mm}
\VUpk{10.2pt}{40}{Vendre i comprar.}
\VUsaut{3.0mm}

%%K t15-04-1 p
4. --- Al mercat l'animació és gran, ja que ha arribat l'hora de la
\VUgras{subhasta} \textsuperscript{(24)}. Les \VUgras{mestresses de casa}
\textsuperscript{(26)} que vénen aquí a fer la compra s'aturen de pas davant de la parada de
la \VUgras{tripaire} \textsuperscript{(25)} per comprar \VUgras{tripes} o un
\VUgras{cap de vedella}.

%%K t15-05-1 p
5. --- Un \VUgras{cuiner} \textsuperscript{(27)} gros regateja una \VUgras{tonyina} bellíssima
amb la \VUgras{peixatera} \textsuperscript{(28)}, que apuja els preus al seu grat. Ha rebut
una gran partida d'\VUgras{anguiles, llenguados, truites} i \VUgras{carpes}. El seu
\VUgras{bacallà} i els seus
\VUgras{musclos} són molt frescos.

%%K t15-tit-3 sub
\VUsaut{4.0mm}
\VUpk{10.2pt}{40}{El barat i el car.}
\VUsaut{3.0mm}

%%K t15-06-1 p
6. --- La \VUgras{pollastrera} \textsuperscript{(29)}, instal·lada al seu costat, és molt
honrada. Quan ven un \VUgras{gall dindi}, una \VUgras{oca}, un \VUgras{ànec} o un simple
\VUgras{pollastre}, no demana més que el preu just.

%%K t15-07-1 p
7. --- A la cantonada de la plaça, al primer pis, una \VUgras{llenceria}
\textsuperscript{(30)} s'hi ha establert fa poc, on els compradors estan sempre segurs de
trobar un assortiment molt bo de \VUgras{estovalles},
\VUgras{tovallons, davantals} i \VUgras{draps de cuina}. Al mateix pis, un
\VUgras{ganiveter} \textsuperscript{(31)} ha muntat el seu taller. Esmola els
\VUgras{ganivets} de les verdulaires i fa als hortolans \VUgras{podalls} de
l'\VUgras{acer} més dur.

%%K t15-tit-4 sub
\VUsaut{4.0mm}
\VUpk{10.2pt}{40}{La botiga de queviures.}
\VUsaut{3.0mm}

%%K t15-08-1 p
8. --- Sota el seu taller s'ha obert no fa gaire una gran botiga de queviures. Ja no és la
\VUgras{botiga d'ultramarins} \textsuperscript{(32)} senzilla i modesta d'altres temps, que
només venia gèneres corrents --- \VUgras{te} \textsuperscript{(33)}, \VUgras{xocolata}
\textsuperscript{(34)}, \VUgras{pa de sucre} \textsuperscript{(37)} i \VUgras{sabó}
\textsuperscript{(38)}. Aquí es ven de tot: \VUgras{galetes cruixents}, \VUgras{conserves}
\textsuperscript{(35)},
\VUgras{licors} \textsuperscript{(36)}, \VUgras{rom, conyac, kirsch, anís} i
\VUgras{curaçao}. S'hi troba caça amb la mateixa facilitat --- \VUgras{llebre}
\textsuperscript{(39)}, \VUgras{tords} \textsuperscript{(40)}, \VUgras{perdius}
\textsuperscript{(41)}, \VUgras{guatlles} \textsuperscript{(42)}, \VUgras{ànecs salvatges}
\textsuperscript{(43)} o \VUgras{faisà} \textsuperscript{(44)} --- que fruita tropical com
\VUgras{plàtans} \textsuperscript{(46)} i \VUgras{pinyes}.

%%K t15-09-1 p
9. --- Un \VUgras{dependent} \textsuperscript{(45)} molt servicial està disposat a lliurar el
que se li demani: \VUgras{melmelada} \textsuperscript{(47)}, de groselles o de codony,
\VUgras{llard} \textsuperscript{(48)}, \VUgras{cogombrets} \textsuperscript{(49)} o
\VUgras{sardines} \textsuperscript{(50)} en salaó.

%%K t15-09-2 p
Això és molt còmode per a les \VUgras{clientes} \textsuperscript{(51)}, que troben, al costat
dels gèneres més corrents, les primícies més rares i la fruita més saborosa: \VUgras{peres}
\textsuperscript{(52)} sucoses, \VUgras{prunes} \textsuperscript{(53)}, \VUgras{raïm}
\textsuperscript{(54)},
\VUgras{préssecs} \textsuperscript{(55)}, \VUgras{ametlles}
\textsuperscript{(56)} i \VUgras{nous} \textsuperscript{(57)}.

%%K t15-tit-5 sub
\VUsaut{4.0mm}
\VUpk{10.2pt}{40}{Els clients.}
\VUsaut{3.0mm}

%%K t15-10-1 p
10. --- L'\VUgras{arròs} \textsuperscript{(58)} i les \VUgras{llenties} \textsuperscript{(59)}
són vora la caixa d'\VUgras{arengades} \textsuperscript{(60)} fumades; la \VUgras{patata}
\textsuperscript{(61)} i la
\VUgras{mongeta} \textsuperscript{(63)} es toquen amb les \VUgras{pomes}
\textsuperscript{(62)}, les \VUgras{figues} \textsuperscript{(64)}, les
\VUgras{prunes seques} \textsuperscript{(65)} i les \VUgras{avellanes}
\textsuperscript{(66)}. També es troben en aquesta botiga \VUgras{ous} \textsuperscript{(67)}
frescos i \VUgras{olives} \textsuperscript{(68)}, de manera que els \VUgras{cistells}
\textsuperscript{(70)} i les \VUgras{xarxes de la compra} \textsuperscript{(73)} s'omplen a
tota pressa.

%%K t15-11-1 p
11. --- El \VUgras{cafè} \textsuperscript{(72)} d'aquesta excel·lent casa ha guanyat una fama
ben merescuda entre les \VUgras{criades} \textsuperscript{(69)} i les cuineres, perquè el vell
\VUgras{adroguer} \textsuperscript{(71)} el torra ell mateix amb la major cura.

%%K t15-tit-6 sub
\VUsaut{4.0mm}
\VUpk{10.2pt}{40}{Coses de veure.}
\VUsaut{3.0mm}

%%K t15-12-1 p
12. --- No tothom ve al mercat a comprar provisions. En el pintoresc anar i venir del carreró
que porta a la nau, es veu la gent més variada. Vora un bonàs
\VUgras{escorxador} \textsuperscript{(75)} que empeny davant seu un
\VUgras{carretó} \textsuperscript{(74)}, passen dos soldats de permís, molt
orgullosos de lluir els seus vistosos uniformes: l'\VUgras{hússar} \textsuperscript{(76)} amb
el seu \VUgras{dolman} blau i el seu \VUgras{xacó de plomall}, i el \VUgras{caporal de zuaus}
amb els seus calçons amples i un
\VUgras{fes} al cap.

%%K t15-13-1 p
13. --- El \VUgras{forner} \textsuperscript{(80)}, dret al llindar del
\VUgras{forn} \textsuperscript{(78)}, acaba de pastar la pasta del seu \VUgras{pa}
\textsuperscript{(79)} i de treure del forn els \VUgras{croissants} \textsuperscript{(81)},
els \VUgras{panets} \textsuperscript{(82)} i les
\VUgras{fogasses} \textsuperscript{(83)} que són a l'aparador.

%%K t15-14-1 p
14. --- Damunt del forn viu una hàbil \VUgras{cosidora} \textsuperscript{(84)} que empra
diverses \VUgras{brodadores} \textsuperscript{(85)}. Al seu costat viu una
\VUgras{bugadera i planxadora} \textsuperscript{(86)}, que emmidona la \VUgras{roba
blanca} \textsuperscript{(88)} després de rentar-la i eixugar-la, i la planxa amb una
\VUgras{planxa} \textsuperscript{(87)}.

%%K t15-tit-7 sub
\VUsaut{4.0mm}
\VUpk{10.2pt}{40}{Carn, peix i verdures.}
\VUsaut{3.0mm}

%%K t15-15-1 p
15. --- A la \VUgras{carnisseria} \textsuperscript{(89)} de la planta baixa es ven carn de
tota mena --- \VUgras{be} \textsuperscript{(90)}, \VUgras{vedella} \textsuperscript{(94)} o
\VUgras{vedell} \textsuperscript{(92)} --- en
\VUgras{cuixes de xai, filets, rostits} i \VUgras{costelles}. El \VUgras{mosso de
carnisseria} \textsuperscript{(93)} talla tota aquesta carn i aparta els
\VUgras{ossos de moll} \textsuperscript{(96)}. A la porta de la carnisseria hi ha
una \VUgras{ostrera} \textsuperscript{(97)} que ven \VUgras{ostres} \textsuperscript{(98)}.
Les obre si se li demana.

%%K t15-16-1 p
16. --- Tots aquests comerciants paguen una patent o un dret de parada; per això el
\VUgras{guàrdia} \textsuperscript{(100)} persegueix sense pietat les pobres
\VUgras{verdulaires ambulants} \textsuperscript{(99)}, que els fan la competència
portant en els seus carretons \VUgras{pastanagues, raves, naps, porros} o
\VUgras{manats d'espàrrecs, pèsols frescos, espinacs, agrelles, créixens} i
\VUgras{julivert}.

%%K t15-tit-8 sub
\VUsaut{4.0mm}
\VUpk{10.2pt}{40}{Compte amb el guàrdia!}
\VUsaut{3.0mm}

%%K t15-17-1 p
17. --- El noi que ven \VUgras{alls} \textsuperscript{(101)} i \VUgras{cebes} sap molt bé que
a ell tampoc no li està permès de vendre la seva mercaderia vora el mercat.

%%K t15-17-1 p suite
Es posa a córrer, a risc de tombar el cistell de \VUgras{crancs},
\VUgras{escamarlans} i \VUgras{gambetes} de la \VUgras{venedora}
\textsuperscript{(102)} de \VUgras{llagostes} i \VUgras{llamàntols}.

%%K t15-18-1 p
18. --- Tothom s'afanya i arma un gran aldarull; però això no impedeix de sentir clarament el
reclam del \VUgras{vidrier} \textsuperscript{(103)} ambulant, ni el pregó del
\VUgras{carboner} \textsuperscript{(104)}, que acaba de deixar un moment el seu carro per
lliurar un \VUgras{sac de carbó} \textsuperscript{(105)}.
\VUsaut{6.0mm}
\VUfilet{22mm}
